\documentclass[12pt,a4paper]{article}
\usepackage[margin=2.15cm]{geometry}
\usepackage{fontspec}
\setmainfont{Liberation Serif}
\setmonofont[Scale=0.89]{DejaVu Sans Mono}
\usepackage{setspace}
\setstretch{1.06}
\usepackage{amsmath}
\usepackage{booktabs}
\usepackage{array}
\usepackage{longtable}
\usepackage{hyperref}
\usepackage{listings}
\usepackage{xcolor}
\usepackage{fancyhdr}
\usepackage{enumitem}
\usepackage{titlesec}
\usepackage{caption}
\usepackage{float}
\usepackage{needspace}
\hypersetup{hidelinks, pdfauthor={Ermicev Sofia}, pdftitle={Operating Systems - Laboratory Work No. 18}}
\definecolor{codebg}{RGB}{245,245,245}
\definecolor{codeframe}{RGB}{225,225,225}
\lstdefinestyle{terminal}{
 basicstyle=\ttfamily\footnotesize,
 backgroundcolor=\color{codebg},
 frame=single,rulecolor=\color{codeframe},
 breaklines=true,columns=fullflexible,keepspaces=true,
 showstringspaces=false,upquote=true,aboveskip=7pt,belowskip=7pt,
 xleftmargin=2pt,xrightmargin=2pt
}
\lstdefinestyle{source}{
 basicstyle=\ttfamily\footnotesize,
 backgroundcolor=\color{codebg},frame=single,rulecolor=\color{codeframe},
 breaklines=true,columns=fullflexible,keepspaces=true,showstringspaces=false,
 numbers=left,numberstyle=\tiny\color{gray},numbersep=9pt,
 xleftmargin=17pt,aboveskip=8pt,belowskip=8pt
}
\captionsetup{font=small}
\setlength{\parindent}{1.1cm}
\setlength{\parskip}{4pt}
\titleformat{\section}{\normalfont\Large\bfseries}{\thesection.}{0.6em}{}
\titleformat{\subsection}{\normalfont\large\bfseries}{\thesubsection.}{0.6em}{}
\pagestyle{fancy}
\fancyhf{}
\fancyfoot[C]{\thepage}
\renewcommand{\headrulewidth}{0pt}


\begin{document}
\begin{titlepage}
\centering
\vspace*{0.55cm}
{\large Ministerul Educației și Cercetării al Republicii Moldova\par}
\vspace{0.2cm}
{\large Universitatea Tehnică a Moldovei\par}
\vspace{0.2cm}
{\large Facultatea Calculatoare, Informatică și Microelectronică\par}
\vspace{0.2cm}
{\large Program: Software Engineering\par}
\vspace{1.55cm}
{\Large\bfseries Operating Systems\par}
\vspace{0.65cm}
{\LARGE\bfseries Laboratory Work No. 18\par}
\vspace{0.35cm}
{\Large\bfseries Build a User-Level Thread Switch\par}
\vspace{0.20cm}
{\large Track B: Build the OS (xv6)\par}
\vfill
\begin{flushright}
\begin{minipage}{0.38\textwidth}
\raggedright
\textbf{Elaborated by:}\\
st. gr. FAF-241\\
\textbf{Ermicev Sofia}
\vspace{0.5cm}

\textbf{Verified by:}\\
\textbf{Patricia Reițman}
\end{minipage}
\end{flushright}
\vfill
{\large Chișinău -- 2026\par}
\end{titlepage}
\tableofcontents
\newpage

\section{Introduction}
The purpose of this laboratory work is to implement a simple \emph{user-level thread library} in xv6, where switching between threads occurs entirely in user space. The program saves and restores the RISC-V registers that must survive function calls, gives each thread an independent stack, and uses cooperative round-robin scheduling. No new xv6 kernel functionality is required for this exercise.

The program was built in the existing \texttt{~/xv6-riscv} source tree with the RISC-V compiler and executed under QEMU. The test run presented below is the actual output recorded from the terminal.

\section{Part 1. Building the Thread Switch}
\subsection{Saved context and thread structure}
Each thread stores a \texttt{struct context} with 14 64-bit registers: the return address \texttt{ra}, stack pointer \texttt{sp}, and \texttt{s0} through \texttt{s11}. These are the RISC-V registers that must be preserved by the called function. Each thread also has its own 8,192-byte, 16-byte-aligned stack and one of the states \texttt{FREE}, \texttt{RUNNING}, or \texttt{RUNNABLE}.

\begin{lstlisting}[style=source,language=C]
struct context {
    uint64 ra, sp;
    uint64 s0, s1, s2, s3, s4, s5;
    uint64 s6, s7, s8, s9, s10, s11;
};

struct thread {
    struct context ctx;
    char stack[8192] __attribute__((aligned(16)));
    int state;
};
\end{lstlisting}

\subsection{Thread creation}
The \texttt{thread\_create()} function finds a free slot and initializes the context of a new thread. At its first context switch, the restored \texttt{ra} points to the thread's entry function, while \texttt{sp} points to the top of its own stack.

\begin{lstlisting}[style=source,language=C]
t->state = RUNNABLE;
t->ctx.ra = (uint64)func;
t->ctx.sp = (uint64)(t->stack + STACK_SIZE);
\end{lstlisting}

\subsection{Cooperative round-robin scheduling}
The function \texttt{thread\_schedule()} searches the thread array in round-robin order, starting after the current thread. When it finds a runnable thread, it marks that thread \texttt{RUNNING}, updates \texttt{current\_thread}, and invokes the context switch. A yielding thread first changes its own state back to \texttt{RUNNABLE}.

\begin{lstlisting}[style=source,language=C]
if (current_thread != next_thread) {
    struct thread *old = current_thread;
    next_thread->state = RUNNING;
    current_thread = next_thread;
    thread_switch(&old->ctx, &next_thread->ctx);
}
\end{lstlisting}

\subsection{RISC-V context switch in Assembly}
The \texttt{thread\_switch} procedure accepts pointers to the old and new contexts in registers \texttt{a0} and \texttt{a1}. It stores the old thread's registers using \texttt{sd} (store doubleword), reloads the next thread's registers using \texttt{ld} (load doubleword), and executes \texttt{ret}. Since each context field is eight bytes, the offset for \texttt{s11} is 104 bytes.

\lstinputlisting[style=source]{uthread_switch.S}

\subsection{Makefile integration}
The existing xv6 Makefile was extended, preserving earlier laboratory programs. Its \texttt{UPROGS} list includes \texttt{\$U/\_uthread}, and the custom rule links \texttt{uthread.o}, \texttt{uthread\_switch.o}, and the standard xv6 user library.

\begin{lstlisting}[style=terminal]
$U/_uthread: $U/uthread.o $U/uthread_switch.o $(ULIB) $U/user.ld
    $(LD) $(LDFLAGS) -T $U/user.ld -o $U/_uthread \
      $U/uthread.o $U/uthread_switch.o $(ULIB)
    $(OBJDUMP) -S $U/_uthread > $U/uthread.asm
\end{lstlisting}

\section{Part 2. Running the Threads}
\subsection{Compilation and xv6 startup}
The standalone binary was built, then the whole xv6 file-system image was generated and QEMU launched:

\begin{lstlisting}[style=terminal]
$ make user/_uthread
riscv64-linux-gnu-gcc ... -c -o user/uthread.o user/uthread.c
riscv64-linux-gnu-gcc -c -o user/uthread_switch.o user/uthread_switch.S
riscv64-linux-gnu-ld ... -o user/_uthread ...
riscv64-linux-gnu-objdump -S user/_uthread > user/uthread.asm

$ make qemu
xv6 kernel is booting
hart 2 starting
hart 1 starting
init: starting sh
$ uthread
\end{lstlisting}

The commands completed successfully; the displayed build output is abbreviated to keep the report readable.

\subsection{Actual execution results}
Three new threads announced their start. Because the functions call \texttt{thread\_yield()} explicitly, their output interleaved in a repeatable round-robin sequence:

\begin{lstlisting}[style=terminal]
thread_a started
thread_b started
thread_c started
thread_c 0
thread_a 0
thread_b 0
thread_c 1
thread_a 1
thread_b 1
thread_c 2
thread_a 2
thread_b 2
...
thread_c 98
thread_a 98
thread_b 98
thread_c 99
thread_a 99
thread_b 99
thread_c: exit after 100
thread_a: exit after 100
thread_b: exit after 100
thread_schedule: no runnable threads
$
\end{lstlisting}

Each thread printed the values 0--99 and incremented its own counter 100 times. The final scheduler message appears because all three worker threads marked themselves \texttt{FREE}, leaving no runnable threads. This message matches the expected output in the assignment.

\subsection{Observe: saved registers and the role of \texttt{ra}}
\noindent\textbf{Question.} What exactly does \texttt{thread\_switch} save and load? Why is \texttt{ra} the register that makes execution continue in the other thread?

\medskip
\noindent\textbf{Answer.} \texttt{thread\_switch} saves the current thread's 14 callee-saved RISC-V registers (\texttt{ra}, \texttt{sp}, and \texttt{s0}--\texttt{s11}) into its context structure, then loads those same registers from the next thread's saved context. The restored \texttt{sp} selects the next thread's stack. The restored \texttt{ra} controls the return destination: \texttt{ret} jumps to the address stored in \texttt{ra}. On the first activation, \texttt{thread\_create} sets \texttt{ra} to the thread function, so the initial \texttt{ret} begins that function. On later activations, \texttt{ra} holds the continuation after the previous call to \texttt{thread\_switch}, so execution resumes exactly where the thread yielded.

\section{Conclusion}
This laboratory demonstrated cooperative user-level threading in xv6 without changes to the kernel. The register-saving Assembly routine, per-thread stack initialization, and round-robin scheduler allowed three threads to alternate and complete 100 iterations each. The observed execution agreed with the lab's expected output and illustrated how saving \texttt{ra} and \texttt{sp} enables control to resume in a different thread.

\end{document}
